% BEGIN VOL08-EXPANSION VIII03
\section{Supplementary worked examples}
\label{sec:viii03-pedagogy}

\begin{example}[Power maps on the circle]\label{ex:viii03-ped-01}
The map $f_k:S^1\to S^1$, $f_k(z)=z^k$, has degree $k$.  Positive $k$ wraps counterclockwise $k$ times, negative $k$ reverses orientation and wraps $|k|$ times, and $k=0$ is constant.
\end{example}

\begin{example}[Degree of the antipodal map]\label{ex:viii03-ped-02}
On $S^n$, the antipodal map $A(x)=-x$ has degree $(-1)^{n+1}$.  This follows from viewing $A$ as the restriction of the linear map $-I$ on $\mathbb R^{n+1}$, whose determinant is $(-1)^{n+1}$.
\end{example}

\begin{example}[Composition multiplies degree]\label{ex:viii03-ped-03}
For maps $S^n\xrightarrow{f}S^n\xrightarrow{g}S^n$, induced maps on top homology are multiplication by $\deg f$ and $\deg g$, hence $\deg(g\circ f)=\deg g\,\deg f$.
\end{example}

\section{Graded supplementary exercises}
The following exercises progress from direct computation and construction to proof, hypothesis testing, application, and synthesis.

\subsection*{Standard computations and constructions}

\begin{exercise}[Identity degree]\label{exr:viii03-ped-01}
Compute the degree of $\operatorname{id}_{S^n}$.
\end{exercise}
\begin{hint}
Look at its action on the fundamental class.
\end{hint}
\begin{solution}
The identity induces the identity on $H_n(S^n)\cong\mathbb Z$, so its degree is $1$.
\end{solution}

\begin{exercise}[Constant degree]\label{exr:viii03-ped-02}
Compute the degree of a constant map $S^n\to S^n$ for $n>0$.
\end{exercise}
\begin{hint}
It factors through a point.
\end{hint}
\begin{solution}
Its induced map on $H_n$ is zero, so the degree is $0$.
\end{solution}

\begin{exercise}[Circle power map]\label{exr:viii03-ped-03}
Compute $\deg(z\mapsto z^{-3})$ on $S^1$.
\end{exercise}
\begin{hint}
Count signed windings.
\end{hint}
\begin{solution}
The map winds three times with reversed orientation, so the degree is $-3$.
\end{solution}

\begin{exercise}[Composite]\label{exr:viii03-ped-04}
If $\deg f=-2$ and $\deg g=5$, compute $\deg(g\circ f)$.
\end{exercise}
\begin{hint}
Use multiplicativity.
\end{hint}
\begin{solution}
The degree is $5(-2)=-10$.
\end{solution}

\begin{exercise}[Antipodal parity]\label{exr:viii03-ped-05}
Compute the degree of the antipodal map on $S^4$.
\end{exercise}
\begin{hint}
Use $(-1)^{n+1}$.
\end{hint}
\begin{solution}
$(-1)^5=-1$.
\end{solution}

\subsection*{Proofs}

\begin{exercise}[Homotopy invariance]\label{exr:viii03-ped-06}
Prove homotopic maps $S^n\to S^n$ have the same degree.
\end{exercise}
\begin{hint}
Homotopic maps induce the same map on homology.
\end{hint}
\begin{solution}
Since degree is the integer describing the induced map on $H_n(S^n)$, homotopy invariance of homology gives equal degrees.
\end{solution}

\begin{exercise}[Multiplicativity]\label{exr:viii03-ped-07}
Prove $\deg(g\circ f)=\deg g\deg f$.
\end{exercise}
\begin{hint}
Write each induced homomorphism as multiplication by an integer.
\end{hint}
\begin{solution}
Functoriality gives $(g\circ f)_*=g_*f_*$; composing multiplication by the two degrees multiplies the integers.
\end{solution}

\begin{exercise}[Nonzero degree not null]\label{exr:viii03-ped-08}
Prove a map $S^n\to S^n$ with nonzero degree is not null-homotopic.
\end{exercise}
\begin{hint}
A constant map has degree zero.
\end{hint}
\begin{solution}
If it were null-homotopic, homotopy invariance would force its degree to equal that of a constant map, namely zero.
\end{solution}

\begin{exercise}[Homotopy equivalence degree]\label{exr:viii03-ped-09}
Prove a homotopy equivalence $S^n\to S^n$ has degree $\pm1$.
\end{exercise}
\begin{hint}
Use a homotopy inverse and multiplicativity.
\end{hint}
\begin{solution}
If $g$ is a homotopy inverse, $\deg g\deg f=1$ in $\mathbb Z$, so both degrees are units, hence $\pm1$.
\end{solution}

\subsection*{Counterexamples and hypothesis tests}

\begin{exercise}[Degree determines map pointwise]\label{exr:viii03-ped-10}
Test: two maps of equal degree are equal as functions.
\end{exercise}
\begin{hint}
Degree is a homotopy invariant, much coarser than pointwise data.
\end{hint}
\begin{solution}
False.  Many distinct maps have the same degree.
\end{solution}

\begin{exercise}[Degree zero means constant]\label{exr:viii03-ped-11}
Test: degree zero forces a map $S^n\to S^n$ to be constant.
\end{exercise}
\begin{hint}
It only detects a homotopy class obstruction.
\end{hint}
\begin{solution}
False pointwise; degree zero maps can be nonconstant, although in the sphere-to-sphere setting they are null-homotopic under the standard classification theorem.
\end{solution}

\begin{exercise}[Every integer occurs]\label{exr:viii03-ped-12}
Test: only $-1,0,1$ can occur as degrees of maps $S^n\to S^n$.
\end{exercise}
\begin{hint}
Use wrapping maps or suspension constructions.
\end{hint}
\begin{solution}
False.  Every integer occurs as the degree of a suitable sphere map.
\end{solution}

\subsection*{Applications and investigations}

\begin{exercise}[Surjectivity certificate]\label{exr:viii03-ped-13}
Explain why nonzero degree forces a map $S^n\to S^n$ to be surjective.
\end{exercise}
\begin{hint}
A map missing one point factors through a contractible punctured sphere.
\end{hint}
\begin{solution}
If the image misses a point, the map factors through $S^n\setminus\{p\}\simeq\mathbb R^n$, hence is null-homotopic and has degree zero; contradiction.
\end{solution}

\begin{exercise}[Orientation bookkeeping]\label{exr:viii03-ped-14}
Why is degree a useful compact summary of orientation and covering multiplicity?
\end{exercise}
\begin{hint}
Local sheets contribute signed counts.
\end{hint}
\begin{solution}
At regular values, the degree is the sum of local orientation signs, combining multiplicity and orientation into one homotopy-invariant integer.
\end{solution}

\subsection*{Challenge problems}

\begin{exercise}[Fixed-point style obstruction]\label{exr:viii03-ped-15}
Suppose $f:S^n\to S^n$ is homotopic to the antipodal map.  What is its degree?
\end{exercise}
\begin{hint}
Use homotopy invariance.
\end{hint}
\begin{solution}
It has degree $(-1)^{n+1}$, the degree of the antipodal map.
\end{solution}

\begin{exercise}[Degree synthesis]\label{exr:viii03-ped-16}
A map $f:S^n\to S^n$ satisfies $f\circ f\simeq f$ and has nonzero degree. Determine $\deg f$.
\end{exercise}
\begin{hint}
Translate the homotopy relation into an integer equation.
\end{hint}
\begin{solution}
Let $d=\deg f$.  Multiplicativity and homotopy invariance give $d^2=d$.  Since $d\ne0$, $d=1$.
\end{solution}

% END VOL08-EXPANSION VIII03
